\exercisesection{§ 1 的习题}

\begin{enumerate}
\def\labelenumi{\arabic{enumi})}
\tightlist
\item
  \begin{enumerate}
  \def\labelenumii{\alph{enumii})}
  \tightlist
  \item
    设 (W, S) 为 Coxeter 系统，并设 \(s_{1}, \ldots, s_{r}\) 是 S 的元素。置 \(w = s_{1} \ldots s_{r}\)。证明：若 \(l_{\mathbb{S}}(w) < r\)，则存在两个整数 p、q，满足 \(1 \leqslant p < q \leqslant r\)，使得 \(w = s_{1} \ldots s_{p-1}s_{p+1} \ldots s_{q-1}s_{q+1} \ldots s_{r}\)。证明存在严格递增的整数序列 \(j(1), \ldots, j(k)\)，其各项介于 1 与 r 之间，使得 \((s_{j(1)}, \ldots, s_{j(k)})\) 是 w 的一个约化分解。
  \end{enumerate}
\end{enumerate}

\begin{enumerate}
\def\labelenumi{\alph{enumi})}
\setcounter{enumi}{1}
\tightlist
\item
  设 \((\mathbf{W},\mathbf{S})\) 为 Coxeter 系统，且 \(X,Y,Z\) 为 \(S\) 的三个子集。证明
\end{enumerate}

\[
W _ {X} \cap (W _ {Y}. W _ {Z}) = (W _ {X} \cap W _ {Y}). (W _ {X} \cap W _ {Z})
\]

（证明每个元素 \(w \in W_{Y}.W_{Z}\) 都有一个约化分解

\[
(s _ {1}, \ldots , s _ {h}, t _ {1}, \ldots , t _ {k})
\]

其中 \(s_i \in Y\)、\(t_j \in Z\)，并使用第 8 项命题 7 的推论 1）。

证明

\[
W _ {X}. \left(W _ {Y} \cap W _ {Z}\right) = \left(W _ {X}. W _ {Y}\right) \cap \left(W _ {X}. W _ {Z}\right).
\]

\begin{enumerate}
\def\labelenumi{\arabic{enumi})}
\setcounter{enumi}{1}
\item
  设 \((W,S)\) 为 Coxeter 系统，X 为 S 的一个子集。证明：\(W_{X}\) 在 W 中正规，当且仅当 X 的每个元素都与 S-X 的每个元素交换。
\item
  设 \((\mathbf{W},\mathbf{S})\) 为 Coxeter 系统，\(\mathbf{X},\mathbf{Y}\) 为 \(\mathbf{S}\) 的两个子集。令 \(a\in W\)。证明存在唯一的最短元素 \(w\in W_X a W_Y\)，并且每个元素
\end{enumerate}

\[
w ^ {\prime} \in W _ {X} a W _ {Y}
\]

可以写成 \(w' = xwy\) 的形式，其中 \(x \in W_X, y \in W_Y\)，且 \(l(w') = l(x) + l(w) + l(y)\)（取 \(W_X aW_Y\) 中的最短元素并使用习题 1）。称元素 \(w \in W\) 为 \((X,Y)\)-约化的，如果它是双陪集 \(W_X aW_Y\) 中的最短元素。

证明：若 \(w\) 是 \((X, \emptyset)\)-约化的，则有 \(l(xw) = l(x) + l(w)\) 对所有 \(x \in W_X\) 成立，并且 \(W\) 的每个元素都能唯一地写成 \(xw\) 的形式，其中 \(x \in W_{X}\)，而 w 是 \((X, \varnothing)\)-约化的。证明元素 \(w \in W\) 是 \((X, \varnothing)\)-约化的，当且仅当 \(l(xw) > l(w)\) 对所有 \(x \in X\) 成立（将 w 写成 \(yw'\) 的形式，其中 \(y \in W_{X}\)，而 \(w'\) 是 \((X, \varnothing)\)-约化的）。

证明 \(w \in W\) 是 \((X, Y)\)-约化的，当且仅当 w 同时是 \((X, \varnothing)\)-约化的和 \((\varnothing, Y)\)-约化的。

\begin{enumerate}
\def\labelenumi{\arabic{enumi})}
\setcounter{enumi}{3}
\item
  设 \(n\) 为一个整数 \(\geqslant 2\)。对于任意整数 \(i\) 满足 \(1 \leqslant i \leqslant n - 1\)，以 \(s_i\) 表示 \(i\) 与 \(i + 1\) 在集合 \(\{1, 2, \ldots, n\}\) 中的换位，以 \(H_i\) 表示 \(w \in \mathfrak{S}_n\) 且满足 \(w^{-1}(i) < w^{-1}(i + 1)\) 的 w 的集合；置 \(S = \{s_1, \ldots, s_{n-1}\}\)。证明 \((\mathfrak{S}_n, S)\) 是 Coxeter 系统，并且 \(H_i\) 是 \(w \in \mathfrak{S}_n\) 中满足 \(l(w) < l(s_i w)\) 的那些 w 的集合（使用第 7 项命题 6）。
\item
  设 X 是一个非空集合，W 是 X 上的一组置换。给定 X 上等价关系的集合 R、元素 \(x_{0} \in X\) 以及从 R 到 W 的映射 \(\varphi : H \mapsto s_{H}\)。以 \(R_{0}\) 表示满足如下条件的 \(H \in R\) 的集合：\(s_{\mathrm{H}}(x_{0}) \equiv x_{0} \mod \mathrm{H}'\) 对 R 中所有 \(H' \neq H\) 均成立；以 \(S_{0}\) 表示这些 \(s_{H}\)（其中 H 属于 \(R_{0}\)）所成的集合。作如下假设：
\end{enumerate}

\begin{enumerate}
\def\labelenumi{(\roman{enumi})}
\item
  对任意 \(\mathbf{H} \in \Re\)，模 \(\mathbf{H}\) 有两个等价类被 \(s_{\mathrm{H}}\) 置换，且 \(s_{\mathrm{H}}^{2} = 1\)。
\item
  对所有 \(\mathbf{H} \in \Re\) 和所有 \(w \in \mathbf{W}\)，\(w(\mathbf{H})\) 是 \(\mathbf{H}\) 在 \(w\) 下的变换，属于 \(\Re\) 的等价关系，且 \(s_{w(\mathbf{H})} = ws_{\mathbf{H}}w^{-1}\)。
\item
  对任意 \(w \neq 1\)（在 \(W\) 中），\(H \in \Re\) 中使 \(w(x_0) \not\equiv x_0 \mod H\) 的那些 H 所成的集合是有限的，并且与 \(\Re_0\) 相交。
\end{enumerate}

\begin{enumerate}
\def\labelenumi{\alph{enumi})}
\item
  证明 \((\mathbf{W},\mathbf{S}_0)\) 是 Coxeter 系统（使用第 7 项命题 6）。
\item
  证明长度 \(l_{S_0}(w)\) 等于满足下式的 \(H \in \mathfrak{R}\) 的元素个数：
\end{enumerate}

\[
w (x _ {0}) \not \equiv x _ {0} \bmod . H.
\]

\begin{enumerate}
\def\labelenumi{\alph{enumi})}
\setcounter{enumi}{2}
\tightlist
\item
  设 E 是一个有限集，X 是 X 上全序关系的集合。以 W 表示 E 的置换群，它以显然的方式作用于 E。设 i、j 是 E 的不同元素；如果 X 的两个元素 R、R' 模 \(H_{ij}\) 等价，即满足 \(\mathrm{R}(i,j)\) 与 \(\mathrm{R}'(i,j)\)，或者满足 \(\mathrm{R}(j,i)\) 与 \(\mathrm{R}'(j,i)\)，则称它们等价；以 \(s_{ij}\) 表示 i 与 j 的换位。令 R 为形如 \(H_{ij}\) 的等价关系的集合，并令 \(\varphi\) 为从 R 到 W 的映射，由 \(\varphi(\mathrm{H}_{ij}) = s_{ij}\) 定义；最后令 \(x_{0}\) 为 X 的任意元素。证明如此定义的对象满足假设 (i) 至 (iii)，并重新得到习题 4 的结果。
\end{enumerate}

\begin{enumerate}
\def\labelenumi{\arabic{enumi})}
\setcounter{enumi}{5}
\item
  设 E 是含有 6 个元素的集合，F 是 E 上域 \(F_{5}\) 的射影直线结构的集合。以 \(S_{E}\) 表示 E 的置换群；对任意 \(\sigma \in S_{E}\)，以 \(\tilde{\sigma}\) 表示 \(\sigma\) 通过结构传递在 F 上诱导的置换。证明存在从 E 到 F 的双射 u，并且映射 \(\sigma \mapsto u^{-1}\tilde{\sigma}u\) 是 \(S_{E}\) 的外自同构（若 s 是换位，则 \(u^{-1}\tilde{s}u\) 有三个二元轨道）。
\item
  构造一个群 W 及其两个子集 S、S'，使得 (W, S) 与 (W, S') 是同构的 Coxeter 系统，但不存在将 S 变为 S' 的 W 的内自同构（使用习题 4 和习题 6）。
\item
  构造一个群 W 及其两个子集 S、S'，使得 (W, S) 与 (W, S') 是不同构的 Coxeter 系统，其中一个不可约而另一个可约（令 W 为由 \(\{s, s'\}\) 生成的 12 阶二面体群，其中 s 和 \(s'\) 的阶为 2 且 \(s \neq s'\)，并置 \(S = \{s, s'\}\)、\(S' = \{(ss')^3, s', s'(ss')^2\}\)）。
\item
  设 (W, S) 是矩阵 \((m(s, s'))\) 给出的 Coxeter 系统，并令 \(W^{+}\) 为由 W 中长度为偶数的元素组成的子群。令 \(s_{0} \in S\)。置 \(g_{s} = ss_{0}\)。证明族 \((g_{s})_{s \in S - \{s_{0}\}}\)，以及关系 \(g_{s}^{m(s, s_{0})} = 1\)（当 \(m(s, s_{0}) \neq \infty\)）和关系 \((g_{s} g_{s'}^{-1})^{m(s, s')} = 1\)（当 \(s, s' \in S - \{s_{0}\}\) 且 \(m(s, s') \neq \infty\)），构成 \(W^{+}\) 的一个表示。（令 \(H^{+}\) 为上述表示定义的群。证明存在一个平方为恒等的自同构 \(\alpha\)，其定义在 \(H^{+}\) 上，将 \(g_{s}\) 变为 \(g_{s}^{-1}\)（对所有 \(s \in S - \{s_{0}\}\)）。若 \(H_{\alpha}\) 是 \(\{1, -1\}\) 与 \(H^{+}\) 关于 \(\alpha\) 的半直积，定义互为逆的同态 \(H_{\alpha} \to W\) 和 \(W \to H_{\alpha}\)。）证明：若 S 的元素彼此共轭（参见命题 3），则群 \(W^{+}\) 是 W 的换位子群（注意此时元素 \(g_{s}\) 都是换位子）。
\item
  令 \(\mathfrak{U}_n\) 为由置换 \(w \in \mathfrak{S}_n\) 中符号等于 \(+1\) 的那些置换组成的交错群。证明 \(\mathfrak{U}_n\) 是 \(\mathfrak{S}_n\) 的换位子群（使用习题 4 和 9）。对于任意整数 \(i\) 满足 \(1 \leqslant i \leqslant n - 2\)，置 \(u_i = s_i s_{i+1}\)（沿用习题 4 的记号）。证明族 \((u_i)\)，以及关系 \(u_1^3 = 1\)、\(u_i^2 = 1\)（当 \(i \geqslant 2\) 时）、\((u_i u_{i+1})^3 = 1\)（当 \(1 \leqslant i \leqslant n - 3\) 时）和 \(u_i u_j = u_j u_i\)（当 \(1 \leqslant i \leqslant n - 4\) 且 \(i + 2 \leqslant j \leqslant n - 2\) 时）构成群 \(\mathfrak{U}_n\) 的一个表示（使用习题 9）。
\end{enumerate}

*11) 令 \((\mathbf{W},\mathbf{S})\) 为 Coxeter 系统。令 \(\Gamma_{\infty}\) 为顶点集是 \(\mathbf{S}\) 的图；两个顶点 \(s,s'\) 以一条边相连，当且仅当 \(m(s,s')\neq \infty\)。令 \(\mathrm{S}_{\alpha}\) 为 \(\Gamma_{\infty}\) 的连通分支。证明 \(\mathbf{W}\) 可以认作各 \(\mathrm{W}_{\mathrm{S}_{\alpha}}\) 的自由积。特别地，每个元素 \(w\in \mathbf{W}\) 都能唯一地写成乘积 \(w_{1}\ldots w_{h}\)，其中 \(w_{i}\neq 1\)，\(w_{i}\) 属于 \(\mathrm{W}_{\mathrm{S}_{\alpha_i}}\)，且有 \(\alpha_{i}\neq \alpha_{i + 1}\)（\(1\leqslant i\leqslant h - 1\)）；证明 \(w\) 的长度等于各 \(w_i\) 的长度之和。*

\begin{enumerate}
\def\labelenumi{\arabic{enumi})}
\setcounter{enumi}{11}
\tightlist
\item
  设 \((\mathbf{W},\mathbf{S})\) 为 Coxeter 系统，满足 \(m(s,s')\) 在 \(s\neq s'\) 时为偶数，并令 \(\mathbf{X}\) 为 \(\mathbf{S}\) 的一个子集。证明存在唯一的同态 \(f_{\mathbf{X}}\)，从 \(\mathbf{W}\) 到 \(\mathbf{W}_{\mathbf{X}}\)，满足 \(f_{\mathbf{X}}(s) = s\)（当 \(s\in \mathbf{X}\)）且 \(f_{\mathbf{X}}(s) = 1\)（当 \(s\in \mathbf{S} - \mathbf{X}\)）。推出 \(\mathbf{W}\) 是 \(\mathbf{W}_{\mathbf{X}}\) 与 \(f_{\mathbf{X}}\) 的核的半直积。证明：若 \(\mathbf{X}\subset \mathbf{Y}\subset \mathbf{S}\)，则存在唯一的同态 \(f_{\mathbf{X},\mathbf{Y}}\)，从 \(\mathbf{W}_{\mathbf{Y}}\) 到 \(\mathbf{W}_{\mathbf{X}}\)，满足 \(f_{\mathbf{X}} = F_{\mathbf{X},\mathbf{Y}}\circ f_{\mathbf{Y}}\)，并且 \(\mathbf{W}\) 可以认作由 X 遍历 S 的有限子集所得到的滤过集上的 \(W_{X}\) 的射影系统的一个子群。
\end{enumerate}

¶ 13) 设 \((W, S)\) 为 Coxeter 系统。对 \(s, s' \in S\)，按以下规则定义序列 \(\mathbf{a}(s, s')\)：

\begin{enumerate}
\def\labelenumi{(\roman{enumi})}
\item
  若 \(ss'\) 的阶为无穷，则 \(\mathbf{a}(s, s')\) 是空序列；
\item
  若 \(ss'\) 的阶为有限数 m，则序列 \(\mathbf{a}(s, s')\) 的长度为 m，且其偶数编号项（分别为奇数编号项）等于 \(s'\)（分别为 s）。
\end{enumerate}

以 \(a(s, s')\) 表示序列 \(\mathbf{a}(s, s')\) 的乘积。

\begin{enumerate}
\def\labelenumi{\alph{enumi})}
\item
  证明生成集 S 以及关系 \(s^{2}=1\) 和 \(\mathbf{a}(s,s')=\mathbf{a}(s',s)\) 构成群 W 的一个表示。
\item
  令 \(q\) 为一个整数 \(\geqslant 1\)。令 \(\Sigma_q\) 为由 \(q\) 个 \(S\) 的元素组成的序列的集合，并令 \(R_q\) 为 \(\Sigma_q\) 上的最小等价关系，使得形如 \((A, \mathbf{a}(s, s'), B)\) 和 \((A, \mathbf{a}(s', s), B)\) 的序列（其中 \(s, s' \in S\)，且 \(A\) 和 \(B\) 是由 \(S\) 的元素组成的序列）等价。令 \(\Sigma_q^r\) 为序列 \(\mathbf{s} = (s_1, \ldots, s_q)\) 的集合，其中 \(w(\mathbf{s}) = s_1 \ldots s_q\) 的长度为 \(q\)。证明序列 \(\mathbf{s}, \mathbf{s}' \in \Sigma_q^r\) 模 \(R_q\) 等价，当且仅当 \(w(\mathbf{s}) = w(\mathbf{s}')\)（按 \(q\) 归纳，并应用第 5 项命题 5）。
\item
  证明序列 \(\mathbf{s} \in \Sigma_q\) 不属于 \(\Sigma_q^r\)，当且仅当 \(\mathbf{s}\) 模 \(R_q\) 等价于一个含有两个相邻相等项的序列。（按 \(q\) 归纳，并归结为序列 \((s_1, \ldots, s_q)\) 的情形：它不属于 \(\Sigma_q^r\)，但 \((s_1, \ldots, s_{q-1})\) 和 \((s_2, \ldots, s_q)\) 属于 \(\Sigma_{q-1}^r\)；使用习题 1 证明 \(s_1 \ldots s_{q-1} = s_2 \ldots s_q\)，再应用 b)。
\end{enumerate}

*14) 设 \((\mathbf{W},\mathbf{S})\) 为 Coxeter 系统，\((\Gamma ,f)\) 为其 Coxeter 图。令 \(k\) 为一个整数 \(\geqslant 3\)。若 \(a\) 是 \(\Gamma\) 的一条边，则置 \(f_{k}(a) = f(a)\)（当 \(f(a)\neq \infty\) 时），并置 \(f_{k}(a) = k\)（当 \(f(a) = \infty\) 时）。令 \((\mathrm{W}_k,\mathrm{S})\) 为 Coxeter 图等于 \((\Gamma ,f_k)\) 的 Coxeter 系统（第 V 章，§4，第 3 项命题 4 的推论）。证明存在唯一的同态 \(\varphi_{k}\)，从 W 到 \(\mathbf{W}_k\)，在 S 上诱导恒等映射。证明若 \(k\) 整除 \(k^{\prime}\)，则存在唯一的同态 \(\varphi_{k,k^{\prime}}\)，从 \(\mathbf{W}_{k^{\prime}}\) 到 \(\mathbf{W}_k\)，满足 \(\varphi_{k} = \varphi_{k,k^{\prime}}\circ \varphi_{k^{\prime}}\)。证明同态 \((\varphi_{k})\) 从 W 到 \(\mathbf{W}_k\) 的射影极限是单射（使用习题 13）\(c)\)，但一般不是满射（例如在无限二面体群的情形）。*

\begin{enumerate}
\def\labelenumi{\arabic{enumi})}
\setcounter{enumi}{14}
\tightlist
\item
  设 A 是一个集合，C 是 \(\mathfrak{P}(A)\) 的一个子集。C 的元素称为 A 的室，室 C 的子集 F 称为一个面，并将 F 在 C 中的余维定义为 C-F 的基数。若面 F 在 C 中的余维为 1，则称 F 为 C 的一个面板。若两个室 C、C' 有公共面板 F，则称它们相邻：此时或者 \(C = C'\)，或者 \(F = C \cap C'\)。长度为 n 的廊是序列 \(\Gamma = (C_0, C_1, \ldots, C_n)\)，由 \(n + 1\) 个室组成，并满足 \(C_i\) 与 \(C_{i+1}\) 相邻（\(0 \leqslant i \leqslant n - 1\)）。于是称 \(C_0\) 与 \(C_n\) 为 \(\Gamma\) 的端点。廊 \(\Gamma\) 称为单射的，当 \(C_i \neq C_{i+1}\)（\(0 \leqslant i \leqslant n - 1\)）时；若不存在具有相同端点且长度 \textless{} n 的廊，则称其为最短的。
\end{enumerate}

若 A 的每个元素至少属于一个室，且任意两个室都是某个廊的端点，则称 A（连同 C）为一个建筑。两个室 C 与 \(C'\) 之间的距离，是长度 \(d(\mathrm{C},\mathrm{C}')\) 的、端点为 C 与 \(C'\) 的最短廊。

A 的一个子建筑是 A 的子集 D，使得 D 连同 \(\mathfrak{C}\cap\mathfrak{P}(D)\) 构成一个建筑。

\begin{enumerate}
\def\labelenumi{\alph{enumi})}
\item
  证明：若 A 是一个建筑，则一个面在包含它的每个室中的余维相同；因此，我们可以不指明特定的室而谈论 A 的面的余维和面板。建筑 B 到 A 的态射是从 B 到 A 的映射 f，使得 f 在 B 的每个室 C 上的限制是从 C 到 A 的一个室 \(f(\mathrm{C})\) 的双射。证明 B 的一个面在 f 下的像仍是具有相同余维的面。
\item
  若每个面板恰好包含于两个室中，则称建筑 A 为一个公寓。证明：若 A 是一个公寓，则 A 的每个保持 C 不变的自同构（即 A 的每个保持 C 的置换）若固定一个室的所有点，就必为恒等映射。更一般地，设 \(\varphi\) 是 A 的一个自同态，C 是 A 的一个室，且 \(\varphi(a)=a\) 对所有 \(a\in C\) 成立。令 \((\mathrm{C},\mathrm{C}_{1},\ldots,\mathrm{C}_{n})\) 为 A 的一个廊。证明：廊 \((\mathrm{C},\varphi(\mathrm{C}_{1}),\ldots,\varphi(\mathrm{C}_{n}))\) 要么不是单射的，要么 \(\varphi(a)=a\) 对所有属于各 \(C_{i}\) 的并集的 a 都成立。
\end{enumerate}

\begin{enumerate}
\def\labelenumi{\arabic{enumi})}
\setcounter{enumi}{15}
\tightlist
\item
  设 \((W,S)\) 为 Coxeter 系统。对 \(s \in S\)，以 \(W^{(s)}\) 表示 W 的子群 \(W_{S-\{s\}}\)，以 A 表示形如 \(wW^{(s)}\)（其中 \(w \in W\) 且 \(s \in S\)）的 W 的子集所成的集合，并以 C 表示形如 \(C_w = \{wW^{(s)} \mid s \in S\}\)（其中 \(w \in W\)）的 A 的子集所成的集合；我们称这些子集为 A 的室（习题 15）。
\end{enumerate}

\begin{enumerate}
\def\labelenumi{\alph{enumi})}
\item
  证明映射 \(w \mapsto \mathbf{C}_w\) 是从 W 到 \(\mathfrak{C}\) 的双射。
\item
  证明两个不同的室 \(C_{w}\) 与 \(C_{w'}\) 相邻，当且仅当存在 \(s \in S\) 使得 \(w' = ws\)。推出 A（连同 C）是一个公寓（习题 15），我们称其为 (W, S) 的公寓。证明端点为 \(C_{w}\) 与 \(C_{w'}\) 的最短廊的长度等于 \(l_{S}(w^{-1}w')\)。
\item
  令 \(\mathfrak{F}\) 为 A 的面的集合，并令 \(F \in \mathfrak{F}\)。证明存在 S 的唯一子集 X 和 W 的一个元素 \(w \in W\)，使得 \(wW_X = \bigcap_{i \in F} i\)。于是称 F 为 X 型。证明 F 的余维等于 X 的基数。证明映射 \(j: F \mapsto \bigcap_{i \in F} i\) 是从 \(\mathfrak{F}\) 到形如 \(wW_X\)（其中 \(w \in W\)、\(X \subset S\)）的 W 的子集的严格递减双射（关于包含关系）。证明：对满足 \(X \subset Y \subset S\) 的每个 Y，每个 X 型面都包含唯一的 Y 型面。
\item
  令 W 通过左平移作用于 A，并置 \(C = C_{e}\)。证明 W 在 C 上单纯传递地作用 \(^{8}\)。令 \(C_{1}, \ldots, C_{n}\) 为 A 的室。证明以下条件等价：
\end{enumerate}

\begin{enumerate}
\def\labelenumi{(\roman{enumi})}
\item
  序列 \(\Gamma = (\mathrm{C}_0 = \mathrm{C},\mathrm{C}_1,\dots ,\mathrm{C}_n)\) 是一个单射廊；
\item
  存在由 S 的元素组成的序列 \(\mathbf{s} = (s_1, \ldots, s_n)\)，使得 \(\mathrm{C}_j = t_j(\mathrm{C}_{j-1})\) 对 \(1 \leqslant j \leqslant n\) 成立，其中 \(t_j\) 是第 4 项公式 (11) 中定义的 \(\Phi(\mathbf{s})\) 的元素。
\end{enumerate}

证明：若这些条件满足，则序列 s 是唯一的，且 \(C_{n} = s_{1} \ldots s_{n}(C)\)。证明廊 \(\Gamma\) 是最短的，当且仅当序列 \(\mathbf{s}(\Gamma) = \mathbf{s}\) 是 \(w = s_{1} \ldots s_{n}\) 的一个约化分解。

\begin{enumerate}
\def\labelenumi{\alph{enumi})}
\setcounter{enumi}{4}
\tightlist
\item
  令 T 为 S 的共轭元的并集。对 \(t \in T\)，A 中在 t 下不变的点集 \(L_{t}\) 称为 t 所定义的墙。证明 \(L_{t}\) 是面板的并集，并且面板 F 包含于 \(L_{t}\) 当且仅当 \(j(\mathrm{F})\) 形如 \(wW_{\{s\}}\)，其中 \(t = wsw^{-1}\)。推出：对任意面板 F，存在唯一的元素 \(t = t(\mathrm{F}) \in \mathrm{T}\) 使得 \(F \subset L_{t}: L_{t}\) 称为 F 的支撑墙。
\end{enumerate}

证明：若 \(w(\mathrm{L}_t) = \mathrm{L}_t\)（其中 \(w \in \mathrm{W}\)），则 \(w = 1\) 或 \(w = t\)。

\(f)\) 令 \(w', w'' \in W\)。置 \(C' = w'(C), C'' = w''(C)\)，并令 \(\Gamma = (C_0 = C', C_1, \ldots, C_n = C'')\) 为端点为 \(C'\) 和 \(C''\) 的单射廊。令 \(t_j\) 为 \(T\) 中定义面板 \(C_j \cap C_{j-1}\) 的支撑墙的元素（当 \(1 \leqslant j \leqslant n\) 时）。证明序列 \(\Psi(T) = (w'^{-1} t_j w')_{1 \leqslant j \leqslant n}\) 与序列 \(\Phi(s(w'^{-1}(\Gamma))\) 重合。对 \(t \in T\)，令 \(n(\Gamma, t)\) 为 \(w'^{-1} tw'\) 在 \(\Psi(\Gamma)\) 中出现的次数。由第 4 项引理 1 推出，数 \((-1)^{n(\Gamma, t)}\) 只依赖于 \(C'\) 和 \(C''\)，而不依赖于 \(\Gamma\)：记其为 \(\eta(C', C'', t)\)。证明关系 \(\eta(C', C'', t) = 1\) 是 \(C'\) 与 \(C''\) 之间的一个等价关系，并且有两个被 \(t\) 置换的等价类。以 \(\mathfrak{C}^+(t)\) 表示包含 \(C\) 的等价类，以 \(\mathfrak{C}^-(t)\) 表示另一个等价类。

证明：对 \(w \in \mathbf{W}\) 和 \(s \in S\)，室 \(w(\mathbf{C})\) 属于 \(\mathfrak{C}^+(s)\)，当且仅当 \(l(sw) > l(w)\)。

\(g)\) 令 \(\mathrm{A}^{+}(t)\)（分别为 \(\mathrm{A}^{-}(t)\)）为属于 \(\mathfrak{C}^{+}(t)\)（分别为 \(\mathfrak{C}^{-}(t)\)）的室的并集（其中 \(t \in \mathrm{T}\)）。证明 \(\mathrm{A}^{+}(t) \cap \mathrm{A}^{-}(t) = \mathrm{L}_{t}\)。（为证明 \(\mathrm{A}^{+}(t) \cap \mathrm{A}^{-}(t) \subset \mathrm{L}_{t}\)，归结到 \(t \in \mathrm{S}\) 的情形。若 \(a \in \mathrm{A}^{+}(t) \cap \mathrm{A}^{-}(t)\)，置 \(a = w\mathrm{W}^{(s)}\)，其中 \(s \in \mathrm{S}\)，且 \(w\) 是 \((\varnothing, \mathrm{S} - \{s\})\)-约化的（习题 3）。若 \(w(\mathrm{C}) \in \mathfrak{C}^{-}(t)\)，则 \(l(tw) < l(w)\)，且 \(w = ts_{1}, \ldots, s_{q}\)，其中 \(s_{j} \in \mathrm{S}\) 且 \(l(w) = q + 1\)。由于 \(a \in \mathrm{A}^{-}(t)\)，存在 \(x \in \mathrm{W}^{(s)}\)，使得 \(wx(\mathrm{C}) \in \mathfrak{C}^{+}(t)\)。于是

\[
l (t w x) = 1 + l (w x) = 1 + l (w) + l (x),
\]

但由于 \(twx = s_1 \ldots s_q x\)，这是矛盾。因此，\(w(\mathbf{C}) \in \mathfrak{C}^+(t)\) 且 \(l(tw) = 1 + l(w)\)。若 \(x \in \mathrm{W}^{(s)}\) 满足 \(l(twx) < l(wx)\)，则由习题 1 推出 \(twx = wx'\)，其中 \(x' \in \mathrm{W}^{(s)}\)，从而 \(ta = a\) 且 \(a \in \mathrm{L}_t\)。

子集 \(\mathbf{A}^{+}(t)\) 和 \(\mathbf{A}^{-}(t)\) 称为由墙 \(L_{t}\) 确定的 A 的两个半部。若 A 的两个点属于（分别不同时属于）两个半部之一，则称它们位于 \(L_{t}\) 的同侧（分别严格位于两侧）。任一面都包含于由 \(L_{t}\) 确定的某个半部中。若两个面包含于不同的半部中，则称它们位于 \(L_{t}\) 的两侧，或称被 \(L_{t}\) 分隔。

\(h\)）令 \(w \in \mathbf{W}\)。证明 \(l(w)\) 等于分隔 \(\mathbf{C}\) 与 \(w(\mathbf{C})\) 的墙的数目。

\begin{enumerate}
\def\labelenumi{\roman{enumi})}
\tightlist
\item
  证明映射 \(\varphi\) 将半部 \(A^{+}(t)\)（分别将 \(A^{-}(t)\)）映到 \((1,t)\)（分别映到 \((-1,t)\)），是从半部集合 \(\mathfrak{M}\) 的 \(A\) 到集合 \(R = \{1, - 1\} \times T\) 的双射（参见第 4 项）。沿用第 4 项引理 1 的记号，证明 \(\varphi(w(M)) = U_w(\varphi(M))\) 对所有 \(w\in W\) 和 \(M\in \mathfrak{M}\) 成立。
\end{enumerate}

\begin{enumerate}
\def\labelenumi{\arabic{enumi})}
\setcounter{enumi}{16}
\item
  沿用习题 16 的记号，并假设 W 有限。令 H 为 A 的墙的集合。对任意 \(H \in H\)，将 A 中由 H 确定且包含 C 的半部 \(H^{+}\) 与之关联。证明可以对 H 的元素编号，使映射 \(j \mapsto \bigcap_{i \leqslant j} H_{i}^{+}\) 在区间 {[}1, \(\operatorname{Card}(\mathfrak{H})\) {]} 上严格递减。（考虑族 \(\mathfrak{F}\)，它由集合 \(\mathrm{H}^{+}\) 的交按包含关系排序而成，并考虑 \((\mathrm{F}_{0},\ldots,\mathrm{F}_{q})\) 这一 \(\mathfrak{F}\) 中长度最大的严格递减序列。对任意 \(\mathrm{H}\in\mathfrak{H}\)，存在 \(i\)，使得 \(\mathrm{H}^{+}\supset\mathrm{F}_{j}\)（当 \(j\geqslant i\) 时）且 \(\mathrm{H}^{+}\not\supset\mathrm{F}_{i-1}\)：证明 \(\mathrm{F}_{i}=\mathrm{H}^{+}\cap\mathrm{F}_{i-1}\)。）
\item
  设 A 是一个公寓（习题 15）。A 的一个折叠是 A 的一个自同态 \(\pi\)，满足 \(\pi^{2} = \pi\)，并且 A 的每个室都是 \(\pi\) 下 0 个或 2 个室的像。
\end{enumerate}

\begin{enumerate}
\def\labelenumi{\alph{enumi})}
\item
  设 \(\pi\) 是 A 的一个折叠，C 是 A 的一个室，且 \(\pi(C) = C\)。对 C 的任意相邻室 \(C'\)，或者 \(\pi(C') = C'\)，或者 \(\pi(C') = C\)；若 \(a \in C\)，则有 \(\pi(a) = a\)。证明：若 \((C_0 = C, C_1, \ldots, C_n)\) 是一个廊，则要么 \(\pi(C_i) = C_i\) 对所有 \(i\) 都成立，要么 \((C_0, \pi(C_1), \ldots, \pi(C_n))\) 不是最短的（事实上有两个相邻且相等的室）。推出：端点在 \(\pi\) 下不变的任意最短廊都在 \(\pi\) 下不变。若 \((C = C_0, C_1, \ldots, C_n)\) 是最短廊且 \(\pi(C_n) \neq C_n\)，则存在一个指标 \(i\)，满足 \(0 \leqslant i < n\)，使得 \(\pi(C_j) = C_j\) 对 \(0 \leqslant j \leqslant i\) 成立，而 \(\pi(C_j) \neq C_j\) 对 \(i < j \leqslant n\) 成立。
\item
  设 \(C_{1}\) 和 \(C_{2}\) 是两个不同的相邻室，\(\pi, \pi'\) 是 A 的两个折叠。假设 \(\pi(C_{0}) = C_{1}\) 且 \(\pi'(C_{1}) = C_{0}\)。令 C 为一个室，并考虑以下三个条件：
\end{enumerate}

\[
\pi (\mathrm{C}) = \mathrm{C};\tag{1}
\]

\[
d (\mathrm{C}, \mathrm{C} _ {1}) <   d (\mathrm{C}, \mathrm{C} _ {2});\tag{2}
\]

\[
\pi^ {\prime} (\mathrm{C}) \neq \mathrm{C}.\tag{3}
\]

证明 (1) \(\Longrightarrow\) (2) \(\Longrightarrow\) (3)，并推出这三个条件等价。证明 \(\pi\)（分别为 \(\pi'\)）是将 \(C_2\)（分别

\(C_{1}\) 映到 \(C_{1}\)（分别为 \(C_{2}\)）的 A 的唯一折叠（假设 (2) 满足，并令 \((C_{1}, C_{2}', \ldots, C_{n}' = C)\) 为一个最短廊：证明 \(\pi'(C_{j+1}')\) 是异于 \(\pi'(C_{j}')\) 且包含面板 \(\pi'(C_{j}' \cap C_{j+1}')\) 的唯一室。证明 \(\pi(\mathfrak{C})\) 和 \(\pi'(\mathfrak{T})\) 构成 A 的室集 T 的一个划分，并且 \(\pi(a) = \pi'(a) = a\) 对所有 \(a \in \pi(A) \cap \pi'(A)\) 成立。证明从 A 到自身的映射与 \(\pi'\) 在 \(\pi(A)\) 上重合，并与 \(\pi\) 在 \(\pi'(A)\) 上重合，是 A 的一个对合自同构。称其为关于面板 \(C_{1} \cap C_{2}\) 的反射。证明它是 A 中固定 \(C_{1} \cap C_{2}\) 的所有点的唯一非平凡自同构（使用习题 15 b)）。

\begin{enumerate}
\def\labelenumi{\alph{enumi})}
\setcounter{enumi}{2}
\tightlist
\item
  设 A 是 Coxeter 系统 (W,S) 所关联的公寓，并沿用习题 16 的记号。设 \(C_{1}\) 和 \(C_{2}\) 是两个相邻室，令 t 为 T 中满足墙 \(L_{t}\) 是 \(C_{1} \cap C_{2}\) 的支撑墙的元素。令 \(M_{j}\) 为由 \(L_{t}\) 确定且包含 \(C_{j}\) 的 A 的半部（j = 1,2）。证明映射 \(\pi\) 由 \(\pi(a) = a\)（当 \(a \in M_{1}\)）和 \(\pi(a) = t(a)\)（当 \(a \in M_{2}\)）定义，是 A 的一个折叠，满足 \(\pi(C_{2}) = C_{1}\)，并且关于面板 \(C_{1} \cap C_{2}\) 的反射是映射 \(a \mapsto t(a)\)。
\end{enumerate}

\begin{enumerate}
\def\labelenumi{\arabic{enumi})}
\setcounter{enumi}{18}
\tightlist
\item
  设 A 为一个公寓。假设对任意两个不同的相邻室 \(C_{1}\) 和 \(C_{2}\)，A 都存在一个折叠（习题 18）将 \(C_{1}\) 映到 \(C_{2}\)。令 C 为 A 的一个室，\((\mathrm{C}_{i})_{i\in I}\) 为与 C 相邻且异于 C 的室族。以 \(s_{i}\) 表示关于面板 \(C\cap C_{i}\) 的反射（习题 18 b)）。置 \(S=\{s_{i}\mid i\in I\}\)，并以 W 表示由 \(s_{i}\) 生成的 A 的自同构群。
\end{enumerate}

\begin{enumerate}
\def\labelenumi{\alph{enumi})}
\item
  证明：对任意室 \(C'\)，存在 \(w \in W\) 使得 \(C' = w(C)\)（对 \(d(C, C')\) 作归纳）。
\item
  证明 (W,S) 是一个 Coxeter 系统（对 \(i\in I\)，置
\end{enumerate}

\[
\mathrm{P} _ {s _ {i}} = \left\{w \in \mathrm{W} \mid w (\mathrm{C}) \subset \pi_ {i} (\mathrm{A}) \right\},
\]

其中 \(\pi_{i}\) 是将 \(C_{i}\) 映到 C 的折叠，并证明命题 6 的假设得到满足：为证明条件 (C')，注意若 \(w \in P_{s_{i}}\) 且 \(ws_{j} \notin P_{s_{i}}\)，则

\[
w (\mathrm{C}) \cap w s _ {j} (\mathrm{C}) \subset \pi_ {i} (\mathrm{A}) \cap s _ {i} \pi_ {i} (\mathrm{A}).
\]

由于 \(w(\mathbf{C})\) 与 \(ws_{j}(\mathbf{C})\) 相邻，故有 \(s_i = ws_jw^{-1}\)（习题 18 b)）。

\begin{enumerate}
\def\labelenumi{\alph{enumi})}
\setcounter{enumi}{2}
\item
  设 F 是室 C 的一个面。证明 F 在 W 中的稳定子群 \(W_{F}\) 由 \(s_{i} \in S\)（满足 \(F \subset C \cap C_{i}\)）生成（令 \(w \in W_{F}\) 满足 \(l_{S}(w) > 1\)，并令 \(i \in I\) 使得 \(w = s_{i}w'\)，其中 \(l(w') = l(w) - 1\)；由命题 6，\(w' \in P_{s_{i}}\)，因而 \(w(\mathrm{C}) \subset s_{i}\pi_{i}(\mathrm{A})\)，\(F \subset \pi_{i}(\mathrm{A}) \cap s_{i}\pi_{i}(\mathrm{A})\)，且 \(s_{i} \in W_{F}\)）。特别地，\(w(\mathrm{C}) = \mathrm{C}\)，当且仅当 w = 1。
\item
  证明映射 \(a \mapsto W_{\{a\}}\) 是从公寓 A 到与 (W, S) 相关的公寓（习题 16）的同构，并且与 W 的作用相容。
\end{enumerate}

\begin{enumerate}
\def\labelenumi{\arabic{enumi})}
\setcounter{enumi}{19}
\tightlist
\item
  设 A 为一个建筑，S 为一个集合。若给定从 A 到 S 的映射 f，使得对 A 的任意室 C，f 在 C 上的限制是从 C 到 S 的双射，则称 A 由 S 编号。若 F 是 A 的一个面，则称 \(f(F)\) 为 F 的类型。设 A 是一个有编号的建筑。若自同态 \(\varphi\) 满足 a 与 \(\varphi(a)\) 对所有 \(a \in A\) 具有相同的类型，则称其为允许的自同态。
\end{enumerate}

\begin{enumerate}
\def\labelenumi{\alph{enumi})}
\item
  设 \(\varphi\) 是 A 的一个自同态。证明：若 A 存在一个室 C，使得 a 与 \(\varphi(a)\) 对所有 \(a \in C\) 具有相同的类型，则 \(\varphi\) 是允许的。证明：若 A 是一个公寓且 \(\pi\) 是 A 的一个折叠（习题 18），则 \(\pi\) 是允许的自同态。
\item
  若 A 的子集 D 满足：对所有 \(a \in A - D\)，都存在 A 的允许自同态 \(\varphi\)，使得 \(\varphi(x) = x\) 对所有 \(x \in D\) 成立，且 \(\varphi(a) \neq a\)，则称 D 为凸的。证明任意凸集的交仍为凸集，并且对 A 的任意子集 D，存在包含 D 的最小凸子集：称其为 D 的凸包，记为 \(\Gamma(D)\)。
\end{enumerate}

\begin{enumerate}
\def\labelenumi{\arabic{enumi})}
\setcounter{enumi}{20}
\tightlist
\item
  设 \((W,S)\) 为 Coxeter 系统，A 为其相关公寓（参见习题 16，沿用其中的记号）。
\end{enumerate}

\begin{enumerate}
\def\labelenumi{\alph{enumi})}
\item
  证明 A 存在唯一的编号（称为典范编号），使面 F 的类型正是习题 16 c) 中定义的类型。今后总假定 A 配备此编号。
\item
  证明 \(\mathbf{A}\) 的允许自同构正是 \(\mathbf{W}\) 的作用。
\item
  设 D 是 A 的包含至少一个室的子集。证明以下条件等价：
\end{enumerate}

\begin{enumerate}
\def\labelenumi{(\roman{enumi})}
\item
  D 是包含 D 的 A 的各半部（习题 16 g)）的交；
\item
  D 是凸的；
\item
  每当两个面 \(F_{1}\) 和 \(F_{2}\) 包含于 D 中时，\(F_{1} \cup F_{2}\) 的凸包包含于 D 中；
\item
  每当一个室 \(\mathbf{C}_1\) 和一个面 \(\mathbf{F}\) 包含于 \(\mathrm{D}\) 中，且 \((C_1, \ldots, C_n)\) 是满足 \(\mathbf{F} \subset \mathbf{C}_n\) 的最短廊时，有 \(C_i \subset D\) 对 \(1 \leqslant i \leqslant n\) 成立。
\end{enumerate}

（为证明 (iii) \(\Longrightarrow\) (iv)，使用习题 15 \(b\)）。为证明 (iv) \(\Longrightarrow\) (i)，用反证法。令 \(D'\) 为包含 D 的 A 的各半部的交；令 \(a \in D' - D\)，令 \(C_0\) 为包含于 D 的一个室，并令 \((C_0, C_1, \ldots, C_n)\) 为满足 \(a \in C_n\) 的最短廊。于是 \(C_j \subset D'\) 对所有 \(j\) 成立。证明存在一个整数 \(j\)，满足 \(0 \leqslant j < n\)，使得 \(C_j \subset D\) 且 \(C_{j+1} \not\subset D\)。令 M（分别为 \(M'\)）为由面板 \(C_j \cap C_{j+1}\) 的支撑墙确定且包含 \(C_j\)（分别包含 \(C_{j+1}\)）的 A 的半部。证明 D \(\not\subset M\)。令 \(b \in D \cap (A - M)\)，并令 \(\Gamma = (C_j, C_1', \ldots, C_p')\) 为满足 \(b \in C_p'\) 的最短廊。于是 \(C_k' \subset D\) 对 \(1 \leqslant k \leqslant p\) 成立，且 \(C_p' \subset M'\)。令 \(\pi\) 为像为 \(M'\) 的 A 的折叠（习题 18 \(c\)）：于是 \(\pi(C_j) = C_{j+1}\)，且廊 \(\pi(\Gamma)\) 不是单射的（习题 18 \(a\)）。若 \(\Gamma' = (C_{j+1}, C_2', \ldots, C_{p-2}', C_p')\)

是从 \(\pi(\Gamma)\) 中删去两个相邻相等室之一所得的廊，则廊 \((\mathrm{C}_{j}, \mathrm{C}_{j+1}, \mathrm{C}_{2}', \ldots, \mathrm{C}_{p-2}', \mathrm{C}_{p}')\) 由 \(\Gamma\) 的定义是最短的。由 (iv) 推出 \(\mathrm{C}_{j+1} \subset \mathrm{D}\)。矛盾。）

\begin{enumerate}
\def\labelenumi{\arabic{enumi})}
\setcounter{enumi}{21}
\tightlist
\item
  沿用习题 16 和 21 的记号。
\end{enumerate}

\begin{enumerate}
\def\labelenumi{\alph{enumi})}
\item
  设 \(t \in \mathbf{T}\)，\(w \in \mathbf{W}\)。证明室 \(\mathbf{C}\) 与 \(w(\mathbf{C})\) 被墙 \(\mathbf{L}_t\) 分隔，当且仅当 \(l(tw) < l(w)\)（使用由半部 \(\mathbf{A}^+(t)\) 确定的折叠）。
\item
  设 \(w_0 \in W\)。证明以下条件等价：
\end{enumerate}

\[
l (w w _ {0}) = l (w _ {0}) - l (w) \quad \text {   for   all   } w \in \mathrm{W};\tag{i}
\]

\[
l (t w _ {0}) <   l (w _ {0}) \quad \text {   for   all   } t \in \mathrm{T};\tag{ii}
\]

当 \(t \in \mathrm{T}\) 时，室 \(\mathrm{C}\) 与 \(w_0(\mathrm{C})\) 被墙 \(\mathbf{L}_t\) 分隔 (iii)

（使用习题 16 h) 证明 (iii) \(\Longrightarrow\) (i)。）

证明这样的元素 \(w_{0}\) 唯一，并且它存在当且仅当 W 是有限的。此时它是 W 中长度最大的元素，并由下式刻画：

\[
l (s w _ {0}) <   l (w _ {0}) \quad \text {   for   all   } s \in S.\tag{iv}
\]

此外，\(w_0^2 = 1\)、\(w_0Sw_0 = S\)，且 \(l(w_0) = \mathrm{Card}(T)\)。

\begin{enumerate}
\def\labelenumi{(\alph{enumi})}
\setcounter{enumi}{2}
\tightlist
\item
  假设 W 是有限的。证明：对任意室 \(C_{0}\)，存在唯一的室 \(-C_{0}\)，使得 \(C_{0} \cup (-C_{0})\) 不包含于 A 的任何半部。证明存在 A 的唯一对合自同构 \(\varphi\)，使得 \(\varphi(C_{0}) = -C_{0}\)（对任意室 \(C_{0}\)），且对 A 的任意墙 L 有 \(\varphi(L) = L\)。置 \(\varphi(a) = -a\)（对 \(a \in A\)）。若 F 是一个面，则称面 \(-F = \varphi(F)\) 与 F 对置。
\end{enumerate}

\begin{enumerate}
\def\labelenumi{\alph{enumi})}
\setcounter{enumi}{3}
\tightlist
\item
  令 \(C_{0}\) 为 A 的一个室，F 为 \(C_{0}\) 的一个划分。证明 \(C_{0} \cup (-F)\) 的凸包是由墙 L 确定且包含 \(C_{0}\) 的 A 的半部，其中 L 是 F 的支撑墙。
\end{enumerate}

\begin{enumerate}
\def\labelenumi{\arabic{enumi})}
\setcounter{enumi}{22}
\item
  沿用习题 16 的记号。令 \(\operatorname{Aut}(A)\) 为公寓 \(A\) 的自同构群。证明：若 \(\varphi \in \operatorname{Aut}(A)\)，则 \(\varphi\) 置换 \(A\) 的墙，并且 \(\varphi t\varphi^{-1} \in T\) 对所有 \(t \in T\) 成立（使用习题 18）。推出 \(W\)（视为 \(\operatorname{Aut}(A)\) 的一个子群）是正规子群，且 \(\operatorname{Aut}(A)\) 是自同构子群 \(E\)（保持室 \(C\)）与 \(W\) 的半直积。证明 \(\operatorname{Aut}(A)\) 在 \(W\) 上的作用给出从 \(E\) 到 Coxeter 系统 \((W, S)\) 的自同构群的同构，或者给出该 Coxeter 图 \((W, S)\) 的自同构群的同构（另参见习题 19）。
\item
  结构建筑是指配备了一个子建筑集合 U 的建筑 I，且该集合满足以下条件：
\end{enumerate}

(SB 1) 子建筑 \(A \in \mathfrak{U}\) 是公寓。

(SB 2) I 的任意两个室都包含于 U 的至少一个元素中。(SB 3) 每当 \(A_{1}, A_{2} \in U\) 且 \(A_{1} \cap A_{2}\) 包含一个室时，就存在从 \(A_{1}\) 到 \(A_{2}\) 的同构，固定 \(A_{1} \cap A_{2}\) 的各点。

设 \((\mathbf{I},\mathfrak{U})\) 为一个结构建筑。\(\mathfrak{U}\) 的元素称为 \((\mathbf{I},\mathfrak{U})\) 的公寓，或简称为 \(\mathbf{I}\) 的公寓。

\begin{enumerate}
\def\labelenumi{\alph{enumi})}
\item
  证明 I 的任意两个公寓同构。令 C 为 I 的一个室，A 为包含 C 的 I 的一个公寓。证明存在唯一的 I 的自同态 \(\rho\)（称为以 C 为中心从 I 到 A 的收缩），使得 \(\rho(a)=a\) 对所有 \(a\in A\) 成立，并且对任意包含 C 的公寓 \(A'\)，\(\rho\) 在 \(A'\) 上的限制是从 \(A'\) 到 A 的同构（注意，由 (SB 2) 和习题 15 b)，对任意包含 C 的公寓 \(A'\)，存在唯一的同构 \(\rho_{A'}\) 的 \(A'\)（固定 C 的所有点））。证明 \(\rho^{2}=\rho\)，且 \(\rho^{-1}(C)=C\)。
\item
  设 A 为 I 的一个公寓，\(C_{0}\) 为一个室，F 为包含于 A 中的一个面。令 \((C_{0}, C_{1}, \ldots, C_{n})\) 为满足 \(F \subset C_{n}\) 的最短廊。证明 \(C_{i} \subset A\) 对 \(1 \leqslant i \leqslant n\) 成立（用反证法：若 \(C_{i} \subset A\) 而 \(C_{i+1} \notin A\)，则考虑以异于 \(C_{i}\) 且包含 \(C_{i} \cap C_{i+1}\) 的 A 的室为中心从 I 到 A 的收缩）。
\item
  设 A 为 I 的一个公寓，C 为 A 的一个室，F 为 C 的一个面板，C' 为 I 的一个室，\(\Gamma = (C_0, \ldots, C_n = C')\) 为满足 \(F \subset C_0\) 的最短廊。证明以 C 为中心从 I 到 A 的收缩 \(\rho\) 将 \(\Gamma\) 变为一个廊
\end{enumerate}

\[
\left(\mathrm{C} _ {0} ^ {\prime}, \dots , \mathrm{C} _ {n} ^ {\prime} = \rho (\mathrm{C} ^ {\prime})\right)
\]

该廊满足 \(F \subset C_{0}^{\prime}\) 且长度最短（考虑公寓 \(A^{\prime}\)，它包含 \(C^{\prime}\) 和 F，并应用 b) 以及 \(\rho\) 在 \(A^{\prime}\) 上的限制是同构这一事实）。

\begin{enumerate}
\def\labelenumi{\alph{enumi})}
\setcounter{enumi}{3}
\tightlist
\item
  设 A 为一个公寓，\(C_{1}\) 和 \(C_{2}\) 为包含于 A 中的两个不同相邻室，\(C'\) 为包含 \(C_{1} \cap C_{2}\) 且异于 \(C_{1}\) 和 \(C_{2}\) 的一个室，\(A_{i}'\) 为包含 \(C_{i}\) 和 \(C'\) 的一个公寓（i = 1, 2）。令 \(\varphi_{i}\)（分别为 \(\psi_{i}\)）为从 I 到 A（分别到 \(A_{i}'\)）且中心为 \(C_{i}\)（分别为 \(C'\)）的收缩，并令 \(\rho_{i}\) 为 \(\varphi_{i} \circ \psi_{i}\) 在 A 上的限制。令 C 为 A 的一个室；证明：若
\end{enumerate}

\[
d (\mathrm{C}, \mathrm{C} _ {1}) \leqslant d (\mathrm{C}, \mathrm{C} _ {2}),
\]

则 \(\rho_{1}(a) = a\) 对所有 \(a \in \mathbb{C}\) 成立，并且 \(d(\mathbb{C}, \mathbb{C}_{1}) < d(\mathbb{C}, \mathbb{C}_{2})\)（考虑最短廊 \(\Gamma\)，其端点为 \(\mathbb{C}\) 和 \(\mathbb{C}_{1}\)，并应用 \(c\)）以证明 \(\rho_{1}(\Gamma)\) 是最短的；然后使用习题 15 b)）。证明：若 \(d(\mathbb{C}, \mathbb{C}_{2}) < d(\mathbb{C}, \mathbb{C}_{1})\)，则 \(\rho_{1}(\mathbb{C}) \neq \mathbb{C}\) 且 \(\rho_{1}^{2}(\mathbb{C}) = \rho_{1}(\mathbb{C})\)（考虑最短廊 \(\Gamma\)，其端点为 \(\mathbb{C}\) 和 \(\mathbb{C}_{2}\)，并使用 \(c\)）以证明 \(\rho_{1}(\Gamma)\) 是一个长度最短的廊，其一个端点等于 \(\rho_{1}(\mathbb{C})\)，另一个端点包含 \(\mathbb{C}_{1} \cap \mathbb{C}_{2}\)；推出 \(d(\mathbb{C}_{1}, \rho_{1}(\mathbb{C}_{1})) \leqslant d(\mathbb{C}_{2}, \rho_{1}(\mathbb{C}_{1}))\)。

证明 \(\rho_{i}\) 是 A 的一个折叠（习题 19）（证明 \(A = \rho_{1}(A) \cup \rho_{2}(A)\)，并定义 A 的对合自同构 \(\sigma\)：置 \(\sigma(a) = \rho_{2}(a)\) 当 \(a \in \rho_{1}(A)\) 时，置 \(\sigma(a) = \rho_{1}(a)\) 当 \(a \in \rho_{2}(A)\) 时；证明：若 C 是包含于 \(\rho_{1}(A)\) 的一个室，则 \(\rho_{1}^{-1}(C) = \{C, \rho_{2}(C)\}\)）。

\begin{enumerate}
\def\labelenumi{\alph{enumi})}
\setcounter{enumi}{4}
\tightlist
\item
  设 I 为一个宽敞的结构建筑，也就是说，每个面板至少包含于三个室中。证明存在一个在同构意义下唯一的 Coxeter 系统 (W,S)，使得 I 的公寓同构于与 (W,S) 相关的公寓 \(A_{0}\)（使用 d) 和习题 19）。令 A 为 I 的一个公寓，\(\varphi\) 为从 \(A_{0}\) 到 A 的同构。证明 I 存在唯一的、取值于 S 的编号，使得 \(\varphi(a)\) 与 a 的类型相同，对所有 \(a \in A_{0}\)（取 A 的一个室 C，并使用 b) 证明：若 \(A'\) 和 \(A''\) 是包含 C 的 I 的两个公寓，则 \(A'\) 和 \(A''\) 上将 C 上确定的编号延拓所得的编号在 \(A' \cap A''\) 上一致）。
\end{enumerate}

我们称如此得到的 Coxeter 系统 (W,S) 和编号适配于 I。

证明 a) 中引入的收缩是允许的自同态。证明包含于 I 的一个公寓 A 中的 I 的子集在 I 中凸，当且仅当它在 A 中凸（习题 20）。

\(f\)）沿用 \(e\)）的记号，并令 \(\mathfrak{A}\) 为从 \(\mathbf{A}_0\) 到 I 的各个公寓的允许同构的集合。证明：若 \(\varphi, \psi \in \mathfrak{A}\)，且 C 是一个室、F 是 I 的一个面，并且二者包含于 \(\varphi(\mathbf{A}_0) \cap \psi(\mathbf{A}_0)\)，则存在 \(w \in W\)，使得 \(\psi^{-1}(C) = w\varphi^{-1}(C)\) 且 \(\psi^{-1}(F) = w\varphi^{-1}(F)\)（考虑同构 \(\lambda\)，它从 \(\varphi(\mathbf{A}_0)\) 到 \(\psi(\mathbf{A}_0)\)，固定 \(\varphi(\mathbf{A}_0) \cap \psi(\mathbf{A}_0)\) 的各点，并将习题 21 \(b\)）应用于自同构 \(\psi^{-1} \circ \lambda \circ \varphi^{-1}\) 的 \(\mathbf{A}_0\)。

\begin{enumerate}
\def\labelenumi{\arabic{enumi})}
\setcounter{enumi}{24}
\tightlist
\item
  设 I 是一个集合，令 \(\mathfrak{F}\) 为 I 的有限子集的集合。对任意 \(A \in \mathfrak{F}\)，置 \(\varepsilon(A) = (-1)^{\operatorname{Card}(A)}\)。设 \(G\) 是一个以加法记的阿贝尔群，并令 \(\varphi\) 与 \(\psi\) 为从 \(\mathfrak{F}\) 到 \(G\) 的两个映射。证明以下性质等价：
\end{enumerate}

\[
\varphi (\mathrm{A}) = \sum_ {\mathrm{B} \subset \mathrm{A}} \psi (\mathrm{B}) \quad \text {   for   all   } \mathrm{A} \in \mathfrak {F};\tag{i}
\]

\[
\psi (A) = \sum_ {B \subset A} \varepsilon (A - B) \varphi (B) \quad \text {   for   all   } A \in \mathfrak {F}.\tag{ii}
\]

\begin{enumerate}
\def\labelenumi{\arabic{enumi})}
\setcounter{enumi}{25}
\tightlist
\item
  设 \((W, S)\) 为 Coxeter 系统，且 S 有限。对 W 的任意子集 H，以 \(\mathrm{H}(t)\) 表示由下式定义的整系数形式幂级数
\end{enumerate}

\[
\mathrm{H} (t) = \sum_ {w \in \mathrm{H}} t ^ {l (w)}.
\]

\begin{enumerate}
\def\labelenumi{\alph{enumi})}
\tightlist
\item
  假设 \(\operatorname{Card}(S) = 2\)。证明
\end{enumerate}

\[
\mathrm{W} (t) = \left(1 + t - t ^ {m} - t ^ {m + 1}\right) / (1 - t) \quad \text { if   } \mathrm{W} \text {   is   of   finite   order   } 2 m,
\]

\[
\mathrm{W} (t) = (1 + t) / (1 - t) \quad \text { if   } \mathrm{W} \text {   is   infinite. }
\]

\begin{enumerate}
\def\labelenumi{\alph{enumi})}
\setcounter{enumi}{1}
\tightlist
\item
  假设 W 是有限的。令 \(w_0\) 为 W 中长度最大的元素（习题 22），并令 \(m = l(w_0)\)。证明
\end{enumerate}

\[
\mathrm{W} (t) = t ^ {m} \mathrm{W} (t ^ {- 1})
\]

（使用习题 22）。

\begin{enumerate}
\def\labelenumi{\alph{enumi})}
\setcounter{enumi}{2}
\tightlist
\item
  设 X 是 S 的一个子集。以 \(A_{X}\) 表示 W 的 \((X,\varnothing)\) -约化元素的集合（习题 3），以 \(W_{X}\) 表示由 X 生成的 W 的子群。我们知道（习题 3），元素 \(w \in W\) 属于 \(A_{X}\)，当且仅当 \(l(xw) = l(w) + 1\) 对所有 \(x \in X\) 成立；W 的每个元素 \(w \in W\) 都能唯一地写成 w = uv 的形式，其中 \(u \in W_{X}\)、\(v \in A_{X}\)，且此时 \(l(w) = l(u) + l(v)\)。推出公式
\end{enumerate}

\[
\mathrm{W} (t) = \mathrm{W} _ {\mathrm{X}} (t). \mathrm{A} _ {\mathrm{X}} (t).
\]

\(d\)）沿用上述记号，以 \(B_{X}\) 表示这样的 \(w\in A_X\)，使得 \(l(sw) = l(w) - 1\) 对所有 \(s\in S - X\) 成立。证明

\[
A _ {X} (t) = \sum_ {X \subset Y \subset S} B _ {Y} (t).
\]

推出

\[
\mathrm{B} _ {\mathrm{X}} (t) = \sum_ {\mathrm{X} \subset \mathrm{Y} \subset \mathrm{S}} \varepsilon (\mathrm{Y} - \mathrm{X}) \mathrm{A} _ {\mathrm{Y}} (t) \quad \text { with } \quad \varepsilon (\mathrm{Z}) = (- 1) ^ {\operatorname{Card} (\mathrm{Z})}
\]

（使用习题 25）。

\begin{enumerate}
\def\labelenumi{\alph{enumi})}
\setcounter{enumi}{4}
\tightlist
\item
  假设 W 是有限的，并按 b) 定义 m 和 \(w_{0}\)。证明 \(B_{\varnothing} = \{w_{0}\}\)，且
\end{enumerate}

\[
t ^ {m} = \sum_ {\mathrm{Y} \subset \mathrm{S}} \varepsilon (\mathrm{Y}) \frac {\mathrm{W} (t)}{\mathrm{W} _ {\mathrm{Y}} (t)}
\]

（使用 c) 和 d)）。

\begin{enumerate}
\def\labelenumi{\alph{enumi})}
\setcounter{enumi}{5}
\tightlist
\item
  假设 W 是无限的。证明 \(B_{\varnothing} = \varnothing\)，且
\end{enumerate}

\[
0 = \sum_ {Y \subset S} \frac {\varepsilon (Y)}{W _ {Y} (t)}.
\]

\(g\)）证明形式幂级数 \(W(t)\) 是 \(t\) 的有理函数（使用 \(f\)）并对 \(\operatorname{Card}(\mathbf{S})\) 作归纳。证明此有理函数在 \(t\) 为单位根时恒为零，并且 \(1 / W(\infty)\) 是整数。证明 \(\frac{1}{W(t^{-1})} \in \mathbf{Z}[[t]]\)。

